\section{Verification, source audit, and integration notes}
\label{cont26:sec:verification}

\subsection{Which computations carry mathematical force?}

The article uses three distinct kinds of evidence. The asymptotic and
root-perturbation proofs are mathematical arguments with stated parameter
ranges. A finite exact certificate completes two arguments by proving
rational inequalities on an analytically justified covering or at
specified points. Ordinary numerical evaluation is used to diagnose
formulas, to illustrate asymptotic orders, and to propose identities.
Increasing the precision of that last kind of evidence does not change
its logical status.

The compact certificate for Theorem~\ref{cont26:joint:thm:Q} is a
computer-assisted proof. Its 112 intervals cover the full compact
region; monotonicity of the enclosed functions transfers each endpoint
calculation to the entire interval. The endpoints outside that region
are covered by an analytic estimate and reflection. Thus the proof
does not infer a continuous inequality from a finite sample.
The certificate contains short rational lower bounds, and the verifier
recomputes the stronger rational expressions that imply them.

The nine Lerch sign enclosures play a different role. They prove
sign changes at exact positive rational arguments. The pre-existing
global zero bound, recalled with its proof mechanism in
Section~\ref{cont26:sec:lerch-bifurcations}, then proves completeness and
simplicity of the endpoint zero lists. The continuity and no-escape
arguments, rather than numerical continuation, force an interior
multiple root.

Finally, the certificates for $S_6$ and $S_8$ prove only the inequalities
in Proposition~\ref{cont26:prop:s6-s8-proximity}. They are useful mathematical
statements about explicitly defined real numbers, but a rational
interval of positive width containing zero is compatible with both
equality and inequality. No lower bound for a nonzero period
combination has been established that would turn these enclosures
into a proof of either conjecture.

\begin{center}
\small
\begin{tabular}{>{\raggedright\arraybackslash}p{0.28\linewidth}>{\raggedright\arraybackslash}p{0.41\linewidth}>{\raggedright\arraybackslash}p{0.21\linewidth}}
\toprule
Verification item & Scope & Role\\
\midrule
Harmonic coefficient compiler &
Exact symbolic Gaussian moments through $d_3$; pure-gamma
normalization at each order &
Algebraic cross-check\\[0.35em]
Uniform harmonic quadrature &
24 cases, 75 working decimal digits; small, proportional,
transition, large-ratio, and large-shift cases &
Numerical diagnostic\\[0.35em]
Reflected inequality &
112 rational intervals, degree-40 series with bounded tails;
analytic endpoint coverage &
Finite proof certificate\\[0.35em]
Reflected moment quadrature &
Eight proportional and twelve critical-window cases,
65 working digits &
Numerical diagnostic\\[0.35em]
Lerch endpoint signs &
Nine rational points, 1000 terms each, $2^{-160}$ integer grid
and monotone tails &
Finite proof certificate\\[0.35em]
Lerch endpoint law and transitions &
Nine difference-kernel diagnostics and two ordinary
multiple-root candidate solves &
Numerical diagnostic\\[0.35em]
$S_6,S_8$ residuals &
Exact $N=1100$ and $N=1200$ Euler enclosures,
independent elementary-constant bounds &
Proof of proximity\\[0.35em]
$S_8$ discovery and audits &
Frozen PSLQ vector; fresh 500-digit Euler and independent
220-digit Mellin evaluations &
Conjectural evidence\\
\bottomrule
\end{tabular}
\end{center}

\subsection{Diagnostic error orders}

The numerical values below use relative error
$(\text{approximation})/(\text{quadrature})-1$. This convention is
the reciprocal of the ratio in some theorem statements, and is
recorded explicitly in the JSON data. For the harmonic example
$p=h$, $a=1/2$, the first corrected formula gives
\[
\begin{array}{c|ccc}
h&25&100&400\\ \hline
\text{relative error}
&-2.91878\cdot10^{-5}&-1.92450\cdot10^{-6}&-1.21909\cdot10^{-7}.
\end{array}
\]
Multiplying the depth by four reduces the error by approximately
sixteen, consistently with the proved $p^{-2}$ order.
At $h=400$ and $p=h\log h$, the corrected relative error is
approximately $3.35\cdot10^{-8}$.
The 24-case set also includes varying $a$, very small $p/h$, and
large $p/h$. These checks are deliberately wider than a single
compact-ratio test.

For equal reflected exponents the successive error reductions reported
after Corollary~\ref{cont26:joint:cor:diagonal} agree with the $n^{-1}$ and
$n^{-2}$ orders. In the critical window the meaningful comparison is
against the reference scale $L_{n,m}$, not against the extremely small
unnormalized moment. At $n/(m+1)-\log m\simeq0$, that normalized moment
approaches $e^d$, rather than one.

The numerical programs use scaled integrals and logarithms to prevent
the cancellation or underflow inherent in the defining sums.
Nevertheless, the quadrature diagnostics are not interval quadrature.
Some extreme-case diagnostic errors approach the numerical precision
floor, and their last displayed digits carry no proved error guarantee.
These cases are omitted from the plotted error comparisons and from
judging an empirical power of the remainder.

\subsection{Audit of the pinned manuscript}

The source is fixed at revision
\src{cc34f73596336f2466d9754cb0f3635bd2bedade}.
The provenance manifest records the inspected text files and their
content hashes. The audit concentrated on the statements used here
and on nearby claims about conductor jets, signed kernels, reflected
remainders, and Stieltjes zero geometry. It is not a claim to have
independently reproved every chapter of the book.

\paragraph{The status of $S_4$.}
The current source's \src{chapters/04-S4-proof.tex} gives an exact
octahedral certificate for the mixed weight-five identity. Earlier
reports describing that identity as conjectural have been superseded.
Accordingly, this article uses $S_4$ as proved background and does
not count it among its new results.

\paragraph{A documentation update for $S_6$.}
The numerical paragraph following the conjecture in
\src{chapters/04-cyclotomic-quotients.tex} correctly describes its
220-digit Mellin and 300-digit Euler values as diagnostics. However,
the source's \src{VALIDATION.md} additionally records a separate
900-term exact integer/fraction replay whose enclosure bounds the
\emph{integer-vector} residual by $10^{-260}$.
The chapter should mention that later certificate so that the
narrative and validation ledger express the same evidence level.
The proposed replacement paragraph in the accompanying patch keeps
the numerical calculations distinct from that certificate and
explicitly retains the conjecture status. Our
Proposition~\ref{cont26:prop:s6-s8-proximity} adds a new, separately replayable
bound for the \emph{normalized left-minus-right} difference.
The normalization must be named whenever residual magnitudes are
compared.

\paragraph{The expanding harmonic range.}
The closing discussion in
\src{chapters/04-proportional-depth.tex} correctly warns that its
compact-$\alpha$ theorem does not imply a uniform remainder when
$\alpha=p/h$ grows. It asks for one effective formula across that
range. Theorem~\ref{cont26:thm:uniform} now supplies such a formula, and
\eqref{cont26:eq:expanding-range} gives an explicit answer to the stated
question. On integration, the open-problem sentence should point to
this theorem, while the warning about the scope of the earlier
compact-parameter result should remain.

\paragraph{Fixed reflected exponent versus joint growth.}
The fixed-$m$ theorems and their sharp remainder statements explicitly
hold $m$ fixed. We found no basis for deleting that restriction.
The critical-window theorem gives a concrete reason to preserve it:
there are sequences with $m/n\to0$ for which the normalized leading
fixed-$m$ approximation tends to $e^d\ne1$.
This is a limitation of an attempted extension, not a counterexample
to the theorem actually stated in the source.

\paragraph{Eventual Lerch saturation versus small derivative order.}
The source's eventual zero-count theorem takes the parameter
derivative order large. Corollary~\ref{cont26:cor:bif-kone} instead holds
that order equal to one and lets the Lerch parameter approach zero.
These are different limits. A sentence asserting the maximal count
for every positive deformation at every derivative order would be
false, but the inspected theorem does not make that assertion.
The new unfolding theorem provides a useful explicit boundary to
such a generalization.

No additional substantive mathematical error was identified in the
targeted material. The proposed source patch is therefore a correction
to the account of available evidence, together with integration
guidance for questions now partly or fully answered; it does not
announce an unsupported retraction of an existing theorem.

\subsection{Reproduction and mathematical dependencies}

The package has no network requirement for certificate replay.
From its root directory, run
\begin{quote}\ttfamily
python code/replay\_exact.py
\end{quote}
with ordinary Python 3, without optimization that disables assertions.
This reconstructs the 112 reflected inequalities, the nine Lerch
signs, and both $S_6/S_8$ rational-enclosure schemes, and compares the
results with the stored certificates. The exact proof programs use
only the Python standard library. The separate coefficient compiler
uses SymPy; numerical diagnostics use mpmath, and the exploratory
Lerch coordinate scan uses NumPy and SciPy. Matplotlib reproduces the
figures. Tested versions and individual commands are in the README.

The analytic harmonic theorem and the Lerch unfolding and endpoint
expansion are independent of numerical diagnostics.
The reflected unique-saddle result depends on the finite rational
inequality proof; the critical-window theorem has an independent
analytic proof. The forced Lerch transitions depend on the nine sign
certificates and the proved zero bound. The Gaussian certificate
depends on elementary Euler and elementary-constant error bounds,
not on the numerical success of PSLQ. None of these programs
constitutes a proof of an algebraic relation between evaluated
periods unless an exact symbolic relation is separately supplied.

